\section{Контекст агента и протоколы обновления}
\label{sec:model}
Обозначим инструкции задачи через $P$, текущее наблюдение через $o_t$,
а предшествующие сообщения и результаты инструментов через $H_{<t}$.
При сохранении полной истории
\begin{equation}
 C_t^{\mathrm{full}}=[P,H_{<t},o_t].
 \label{eq:ctx-transcript}
\end{equation}
Структурированная память формирует контекст
\begin{equation}
 C_t^{\mathrm{state},k}=[P,\Sigma_t,\tau_k(H_{<t}),o_t],
 \qquad \tau_0(H)=\varnothing,
 \label{eq:ctx-state}
\end{equation}
где $\tau_k$ выбирает доступные недавние записи, а $\tau_\infty(H)=H$.
При $k=\infty$ сохраняются \emph{и} состояние, \emph{и} весь транскрипт;
контроль с полной историей без состояния --- другая политика
(рис.~\ref{fig:architecture}). Доступ к скрытым рассуждениям не предполагается.
% Adapted from the repository's original two-column architecture figure.
\begin{figure}[tbp]
\centering
\resizebox{\linewidth}{!}{%
\begin{tikzpicture}[
font=\small,
box/.style={draw,rounded corners=2pt,minimum width=3.5cm,minimum height=.62cm,align=center,inner sep=3pt},
hist/.style={box,fill=red!8},
state/.style={box,fill=blue!10,draw=blue!60,thick},
plain/.style={box,fill=black!3},
llm/.style={box,fill=black!8,minimum height=.8cm,thick},
grp/.style={draw,dashed,rounded corners=3pt,inner sep=7pt},
arr/.style={-latex,thick}]
\node[font=\bfseries] at (0,4.2) {Full history};
\node[plain] (p1) at (0,3.25) {task instructions $P$};
\node[hist] (h1) at (0,2.4) {retained transcript $H_{<t}$};
\node[plain] (o1) at (0,1.55) {latest observation $o_t$};
\node[grp,fit=(p1)(h1)(o1)] (g1) {};
\node[llm] (m1) at (0,.2) {language model};
\node[plain] (a1) at (0,-1) {execute action $a_t$};
\draw[arr] (g1.south)--(m1.north);
\draw[arr] (m1)--(a1);
\draw[arr,red!60!black] (a1.west)--++(-1,0)|-(h1.west)
 node[pos=.3,left,font=\scriptsize,align=center]{append\\interaction};
\begin{scope}[xshift=7.5cm]
\node[font=\bfseries] at (0,4.2) {State with optional recent history};
\node[plain] (p2) at (0,3.25) {task instructions $P$};
\node[hist] (h2) at (0,2.4) {recent records $\tau_k(H_{<t})$};
\node[state] (s2) at (0,1.55) {structured state $\Sigma_t$};
\node[plain] (o2) at (0,.7) {latest observation $o_t$};
\node[grp,fit=(p2)(h2)(s2)(o2)] (g2) {};
\node[llm] (m2) at (0,-.7) {language model};
\node[plain] (a2) at (0,-1.9) {execute action $a_t$};
\draw[arr] (g2.south)--(m2.north);
\draw[arr] (m2)--(a2);
\draw[arr,blue!60!black] (m2.east)--++(1.2,0)|-(s2.east)
 node[pos=.25,right,font=\scriptsize,align=left]{model update\\to semantic state};
\node[right=3pt of h2,font=\scriptsize] {$k=0$: empty};
\end{scope}
\node[font=\footnotesize,align=center] at (0,-2.75) {Earlier records remain available\\without rewriting them as state.};
\node[font=\footnotesize,align=center] at (7.5,-2.75) {Earlier records outside the window\\must be represented in state to persist.};
\end{tikzpicture}}
\caption{Context policies. State replaces the older transcript, while the retention parameter $k$ controls the recent suffix. The arrows show information flow, not the number of model calls.}
\label{fig:architecture}
\end{figure}


\subsection{Кто записывает состояние и как оно обновляется}
В первом подходе модель поддерживает семантическое состояние $S_t$.
Гибридная память добавляет записи среды исполнения $V_t$:
\begin{equation}
 \Sigma_t=(S_t,V_t),\qquad
 S_{t+1}=\operatorname{Apply}(S_t,\Delta S_t),\qquad
 V_{t+1}=\operatorname{Record}(V_t,a_t,r_t).
 \label{eq:ownership}
\end{equation}
Модель ведёт гипотезы и незавершённую работу; рантайм сохраняет наблюдаемые
события инструментов $r_t$: от ссылок на файлы до изменений, результатов
команд, диагностик, тестов и неудачных попыток. Обновление $S$ не может
изменить $V$, но полного захвата событий это не гарантирует. Недавняя
история может дополнять оба типа памяти.

Отдельный инструмент принимает изменения состояния; inline-протокол
объединяет обновление и действие в ответе $\{\Delta S_t,a_t\}$, следуя
контракту SKILL.state~\citep{skillstate2026}. Алгоритмы~\ref{alg:context}
и~\ref{alg:inline} --- опорные контракты, а не восстановленные трассы.
\begin{algorithm}[H]
\caption{Context selection and state ownership}
\label{alg:context}
\begin{algorithmic}[1]
\Require Instructions $P$, policy $\pi$, retention setting $k$, update protocol $q$, initial observation $o$
\State Initialize history $H$, semantic state $S$, runtime record $V$
\While{the task has not terminated}
  \If{$\pi=\mathrm{FullHistory}$}
    \State $C \gets [P,H,o]$
  \Else
    \State $\Sigma \gets S$ \textbf{or} $(S,V)$ according to state ownership
    \State $C \gets [P,\Sigma,\tau_k(H),o]$
  \EndIf
  \State $(S',a) \gets \operatorname{PolicyStep}(C,S,\pi,q)$
  \State $r \gets \operatorname{Execute}(a)$; $S \gets S'$
  \If{hybrid memory is enabled}
    \State $V \gets \operatorname{Record}(V,a,r)$
  \EndIf
  \State Append a record of $(o,a)$ to $H$; $o \gets r$
\EndWhile
\end{algorithmic}
\end{algorithm}
\begin{algorithm}[H]
\caption{Inline patch--action contract}
\label{alg:inline}
\begin{algorithmic}[1]
\Require Current context $C$ and semantic state $S$
\State $(\Delta S,a) \gets \operatorname{Model}(C)$
\If{$\operatorname{Valid}(\Delta S,S)$}
  \State \Return $(\operatorname{Apply}(S,\Delta S),a)$
\Else
  \State Keep $S$ unchanged and return a validation observation
  \State Do not execute $a$ from the rejected envelope
\EndIf
\end{algorithmic}
\end{algorithm}


\subsection{Границы контекста и стоимость задачи}
Полная история даёт $O(T^2)$ накопленного входа за $T$ действий;
ограниченное состояние и постоянное окно дают $O(T)$. Ни одна оценка
не фиксирует само $T$. При накопленных некешированном входе $U$, чтении
из кеша $K$, выходе $O$ и тарифах $\alpha,\gamma,\beta$
\begin{equation}
 D=\alpha U+\gamma K+\beta O.
 \label{eq:cost-model}
\end{equation}
Ведение состояния может увеличить выход, снизить переиспользование кеша
или изменить $T$, поэтому меньший пиковый промпт не определяет $D$.
